\section{A Program Logic for Adaptive Games}
\label{sec:pyth-toolkit}

SSProve combines high-level algebraic reasoning about packages with a
relational program logic for their low-level procedures~\cite{ssprove}.  Its
central bridge turns per-procedure exact judgments into perfect
indistinguishability of packages.  We retain that two-level architecture but
add a quantitative bridge needed by noise flooding: a bounded conditional KL
cost for one oracle implementation is lifted to an additive-error bound for a
complete adaptive program.  This section presents that contribution
independently of the FHE instantiation.

The development introduces no new programming language and assumes no logic
rules.  It defines semantic judgments over SSProve's existing
semantics, then verifies the soundness of every rule below.
The resulting interface has
three layers: unary Hoare reasoning establishes support invariants; an
additive-error judgment handles conventional game hops; and a Pythagorean
judgment retains conditional KL information until the end of a computation.

\subsection{SSProve syntax and semantics}
\label{sec:ssprove-substrate}

For completeness, we recall the fragment of SSProve on which our logic is
built~\cite{ssprove}.  SSProve embeds pure expressions shallowly in Rocq and
represents effects by a free monad.  Fix a result type \(A\).  Its raw code has
the following constructors:
\begin{equation}
\label{eq:ssprove-syntax}
\begin{split}
c ::= {}& \mathsf{return}(a)
 \mid \mathsf{call}(p,x,\kappa)
 \mid \mathsf{get}(\ell,\kappa)\\
 &{}\mid \mathsf{put}(\ell,v,c)
 \mid \mathsf{sample}(D,\kappa).
\end{split}
\end{equation}
Here \(a:A\).  An operation signature \(p:S\to T\) contains a procedure
identifier as well as its argument and result types; a call has \(x:S\) and
\(\kappa:T\to\mathsf{RawCode}(A)\).  A typed location
\(\ell:T_\ell\) is read into a continuation
\(\kappa:T_\ell\to\mathsf{RawCode}(A)\), or updated with \(v:T_\ell\).
Finally, sampling takes \(D\in\mathsf{SD}(X)\) and
\(\kappa:X\to\mathsf{RawCode}(A)\), where \(\mathsf{SD}(X)\) denotes discrete
subdistributions on \(X\).  Continuations are Rocq functions, so conditionals,
pure local computation, and structurally terminating loops are supplied by
the ambient language rather than by additional constructors.  Monadic bind
recursively pushes a continuation through this syntax.  We use the customary
notations
\begin{equation}
\label{eq:ssprove-notation}
\begin{aligned}
 x\leftarrow c_1;c_2
   &\equiv\mathsf{bind}(c_1,\lambda x.c_2),\\
 x\leftarrow p(a);c
   &\equiv\mathsf{call}(p,a,\lambda x.c),\\
 x\leftarrow\mathsf{get}\,\ell;c
   &\equiv\mathsf{get}(\ell,\lambda x.c),\\
 \mathsf{put}\,\ell:=v;c
   &\equiv\mathsf{put}(\ell,v,c),\\
 x\leftarrow D;c
   &\equiv\mathsf{sample}(D,\lambda x.c).
\end{aligned}
\end{equation}

An interface is a finite set of typed operation signatures.  The checked type
\(\mathsf{Code}_{L,I}(A)\) pairs raw code with a proof that every accessed
location belongs to the footprint \(L\) and every external call belongs to
the import interface \(I\).  A package \(P:I\to E\) is a finite map from the
procedure names in its export interface \(E\) to well-scoped implementations
\(S\to\mathsf{Code}_{L,I}(T)\).  Sequential composition \(P\circ Q\)
links \(P\)'s calls to procedures implemented by \(Q\).  At code level, the
essential clause is
\begin{equation}
\label{eq:ssprove-link}
\begin{aligned}
 &\mathsf{codeLink}(\mathsf{call}(p,a,\kappa),Q)\\
 &\qquad =
 x\leftarrow Q.p(a);
 \mathsf{codeLink}(\kappa(x),Q),
\end{aligned}
\end{equation}
with \(\mathsf{codeLink}\) acting recursively on the remaining constructors.
Parallel composition combines disjoint export interfaces, while the identity
package forwards every call.  Package validity tracks the import, export, and
memory-footprint side conditions used by these operations.

Each location carries a type and a default initial value.  A heap \(h\) maps
locations to values of their declared types; lookup of an unset location
returns its default.  After linking has eliminated external calls, the
denotation of closed code is a state-transforming subdistribution
\begin{equation}
\label{eq:ssprove-denotation}
 \semantic{c}:\mathsf{Heap}\longrightarrow
  \mathsf{SD}(A\times\mathsf{Heap}).
\end{equation}
Writing \(\delta_z\) for the point distribution at \(z\), its defining
equations are
\begin{align}
 \semantic{\mathsf{return}(a)}_h
   &=\delta_{(a,h)},\notag\\
 \semantic{\mathsf{get}(\ell,\kappa)}_h
   &=\semantic{\kappa(h(\ell))}_h,\notag\\
 \semantic{\mathsf{put}(\ell,v,c)}_h
   &=\semantic{c}_{h[\ell\mapsto v]},\notag\\
 \semantic{\mathsf{sample}(D,\kappa)}_h
   &=\sum_{x}D(x)\,\semantic{\kappa(x)}_h.
\label{eq:ssprove-effect-semantics}
\end{align}
The sums are pointwise sums of subdistributions.  In particular, semantic
bind threads both the returned value and the updated heap:
\begin{align}
 \semantic{x\leftarrow c;\kappa(x)}_h
 &=
 \sum_{(a,h')}
   \semantic{c}_h(a,h')\,
   \semantic{\kappa(a)}_{h'}.
\label{eq:ssprove-bind-semantics}
\end{align}
SSProve defines failure by sampling the null subdistribution, and
\(\mathsf{assert}(b)\) as \(\mathsf{return}(())\) when \(b\) holds and failure
otherwise.  Thus failed assertions have the zero subdistribution and their
continuations are not run.  A closed game is evaluated by resolving its
\(\mathsf{Run}\) procedure, starting from the default heap, and projecting
the result component of \Cref{eq:ssprove-denotation}.

\subsection{Completed semantics and three judgments}

The denotation in \Cref{eq:ssprove-denotation} may have weight below one.  For
a subdistribution \(\mu\) on \(X\), define its completion on
\(X\cup\{\bot\}\) by
\[
 \completed{\mu}(x)=\mu(x),\qquad
 \completed{\mu}(\bot)=1-\sum_x\mu(x).
\]
Completion makes assertion failure observable.
Without this, its becomes tedious (and sometimes tricky) to write down couplings for sub-distributions and then estimate the probabilities of post-conditions.

The unary judgment
\[
 \vDash\mathsf{Hoare}\{P\}\ c\ \{Q\}
\]
is the standard Hoare triple over the SSProve code.
It states that whenever the pre-condition $P$ is satisfied,
the post-condition $Q$ would be satisfied with certainty (i.e. with probability $1$).

The additive-error judgment
\[
 \vDash\AEJ\{P\}\ c_L\approx_\delta c_R\ \{Q\}
\]
states that, from any pair of initial configurations satisfying \(P\), the
completed outputs of \(c_L\) and \(c_R\) admit a coupling in which \(Q\) fails
with probability at most \(\delta\).
For equality post-conditions, this is the
ordinary statistical-distance bound.
We prove consequence, sequencing with loss \(\delta_1+\delta_2\),
triangle, projection, and direct-coupling rules.

The Pythagorean judgment
\begin{equation}
\label{eq:pyth-judgment}
 \vDash\Pyth\{P\}\ c_L\approx_s c_R\ \{Q\}
\end{equation}
captures the Micciancio-Walter \cite{mw17} Pythagorean prerequisites for KL divergence.
It uses a nonnegative tuple \(s=(s_1,\ldots,s_k)\).  For every related input pair,
it supplies two joint transcript distributions whose final
marginals are the programs' outputs.  At coordinate \(i\), for
every common prefix, the conditional KL divergence is finite and at most
\(s_i\).
Outputs on both sides additionally satisfy the post-condition \(Q\).
The transcript is semantic rather than a list of syntactic sampling sites: it may record a
whole program fragment as one coordinate, or expose multiple coordinates for
composition.

This judgment deliberately stores more information than an immediate
statistical-distance bound.  Its sequence rule concatenates tuples (i.e. $s\doubleplus s'$) rather
than adding square roots, allowing one nonlinear conversion after the full
adaptive transcript has been constructed.

\subsection{Core proof rules}
\label{sec:logic-rules}

\Cref{fig:logic-rules} shows the rules that drive the generic development.
They are schematic to suppress SSProve typing, full-mass, support, and heap
side conditions; the artifact contains their concrete statements.

\begin{figure*}[t]
\centering
\small
\begin{mathpar}
\inferrule*[right=\textsc{KL-Sample}]
  {\mathsf{finiteKL}(D_L,D_R) \\
   \KL{D_L}{D_R}\leq\epsilon \\
   \mathsf{weight}(D_L)=\mathsf{weight}(D_R)=1}
  {\vDash\Pyth\{P\}\
     x\leftarrow D_L
     \approx_{[\epsilon]}
     x\leftarrow D_R\ \{Q\}}

\inferrule*[right=\textsc{Pyth-Refl}]
  {P\Longrightarrow x_L=x_R\wedge m_L=m_R \\
   0\leq s_i\ \text{for every }i}
  {\vDash\Pyth\{P\}\ c\approx_s c\ \{Q\}}

\inferrule*[right=\textsc{Pyth-Seq}]
  {\vDash\Pyth\{P\}\ c_L\approx_s c_R\ \{M\} \\
   \vDash\Pyth\{M^{=}\}\ k_L\approx_t k_R\ \{Q\}}
  {\vDash\Pyth\{P\}\
     (x\leftarrow c_L;k_L(x))
     \approx_{s\doubleplus t}
     (x\leftarrow c_R;k_R(x))\ \{Q\}}

\inferrule*[right=\textsc{Shared-Prefix}]
  {\vDash\mathsf{Hoare}\{P_0\}\ c\ \{M\} \\
   \vDash\Pyth\{M^{=}\}\ k_L\approx_s k_R\ \{Q\}}
  {\vDash\Pyth\{P\}\
     (x\leftarrow c;k_L(x))
     \approx_{[0]\doubleplus s}
     (x\leftarrow c;k_R(x))\ \{Q\}}

\inferrule*[right=\textsc{Micciancio--Walter}]
  {\vDash\Pyth\{P\}\ c_L\approx_s c_R\ \{Q\}}
  {\vDash\AEJ\{P\}\ c_L
     \approx_{\sqrt{\norm{s}_1/2}}c_R\ \{o_L=o_R\}}
\end{mathpar}
\caption{Representative checked rules.  Here \(M^{=}\) requires equal
intermediate values and heaps satisfying \(M\), and \(\doubleplus\) denotes
tuple concatenation.  The Rocq theorems additionally expose all typing,
support, invariant, and nonnegativity premises.}
\label{fig:logic-rules}
\end{figure*}

\textsc{KL-Sample} embeds a distribution-level KL bound as a program judgment.
\textsc{Pyth-Refl} relates a program to itself at zero cost.
\textsc{Pyth-Seq} is the essential rule: its proof
constructs a joint transcript for semantic bind and proves conditional bounds
for both valid and zero-mass prefixes.  The development also contains a
shared-continuation rule, consequence rules, closed-program variants,
and the additive-error rules mentioned above.

The mixed Hoare/Pythagorean rules have been useful in practice.  Deterministic
prefixes, table operations, other oracle calls, and postprocessing are proved to
preserve an invariant with the unary Hoare logic and inserted as zero-cost
coordinates in the Pythagorean judgment.
Only the sampling fragment whose distribution changes consumes
KL budget.  In \Cref{lem:one-call}, this decomposition produces the tuple
\((0,\epsilon_{\mathsf{nf}},0)\).

\subsection{The Micciancio-Walter Probability Lemma}

We discuss the probabilistic Pythagorean preservation lemma used to justify the Micciancio-Walter rule in
\Cref{fig:logic-rules}.
Let \(\mathcal P,\mathcal Q\) be joint distributions on \(\Omega^k\), and
write \(\mathcal P_i\mid a\) for coordinate \(i\) conditioned on prefix
\(a\in\Omega^{i-1}\). 

\begin{theorem}[Conditional-coordinate preservation]
\label{thm:pyth-probability}
Suppose \(\mathcal P,\mathcal Q\) have weight one and, for every coordinate
\(i\) and prefix \(a\), the conditional KL divergence is finite and
\[
 \KL{\mathcal P_i\mid a}{\mathcal Q_i\mid a}\leq s_i,
 \qquad s_i\geq0.
\]
Then their final-coordinate marginals satisfy
\[
 \TVD(\mathcal{P}_k,
      \mathcal{Q}_k)
 \leq \sqrt{\frac{\sum_{i=1}^{k}s_i}{2}}.
\]
\end{theorem}

At paper level, the proof expands joint mass as prefix mass times conditional
mass and groups the log likelihood ratio by coordinate.  The resulting chain
bound gives \(\KL{\mathcal P}{\mathcal Q}\leq\sum_i s_i\).  Pinsker bounds the
transcript distance, and data processing gives the displayed final-marginal
bound.  The Rocq proof makes explicit absolute continuity, conditional
zero-mass prefixes, summability of positive and negative parts, and exchanges
between finite coordinate sums and countable support sums.

Our \textsc{Micciancio--Walter} rule instantiates this theorem with the
transcripts from \Cref{eq:pyth-judgment}, then uses maximal coupling to
obtain the additive-error equality judgment.
As with all of our inference rules,
we verify it to be sound with respect to SSProve's program semantics.

\subsection{A derived discrete-Gaussian rule}

For center \(\mu\in\bbZ\) and \(\sigma>0\), we denote the discrete Gaussian distribution by
\[
 D_{\bbZ,\mu,\sigma^2}(x)
 =\frac{\exp(-(x-\mu)^2/(2\sigma^2))}
        {\sum_{z\in\bbZ}\exp(-(z-\mu)^2/(2\sigma^2))}.
\]
The development proves normalization, translation invariance of the
normalizer, summability, the centered first moment, full support, and the
equal-variance KL identity in \Cref{eq:dg-kl}.  The calculation expands the
logarithm of the ratio; equal variances cancel both normalizers, and symmetry
of the centered distribution cancels the linear moment.

These facts yield a reusable program rule for vector sampling.  If
\(u,v\in\bbZ^n\), \(\sigma>0\), and every coordinate satisfies
\[
 \frac{(v_i-u_i)^2}{2\sigma^2}\leq\epsilon,
\]
then the $n$-fold product sampler satisfy a Pythagorean
judgment with budget \([n\epsilon]\).
% The theorem includes injectivity of the SSProve encoding and heap/postcondition premises.
The analysis first constructs
an \(n\)-coordinate transcript, proves the scalar bound coordinatewise, and
then transports and aggregates it into the program-level sampling rule.  The
noise-flooding proof instantiates \(\epsilon=1/(2\gamma^2)\), obtaining
\(\epsilon_{\mathsf{nf}}=n/(2\gamma^2)\) without introducing any unverified assumptions.

\subsection{From one oracle call to an adaptive program}
\label{sec:trace-compiler}

The rules above compose a sequence of sampling operations. An oracle program
does not provide such a sequence: the location and arguments of its next call
may depend on previous answers.  Our compiler is the verified adapter
between these views.  

Intuitively speaking, the pen-and-paper analysis would run the adversary to its next oracle call and return
its query together with a suspended execution.
In SSProve, however, a suspended program is a \emph{continuation} inside the program monad: storing it directly would
require serializing arbitrary Rocq functions into SSProve-compatible data types.
Instead, as an simpler alternative, we choose to record execution traces.
In particular, we record the results of other oracle calls, heap reads
and writes, and samples; replaying them recovers the continuation.
The same-package theorem below proves this representation exact.

Write
\(\mathsf{Comp}_q(A;P_s,P_r)\) for the program obtained by exposing the first
\(q\) calls made by \(A\) to a selected operation, resolving those calls with
\(P_s\), and linking the residual program with \(P_r\).  The compiler first
satisfies an exactness theorem.

\begin{theorem}[Same-package compilation]
\label{thm:compile-same}
Subject to package typing and memory-separation conditions,
\[
 \mathsf{Comp}_q(A;P,P)
 \quad\text{and}\quad
 \mathsf{Link}(A,P)
\]
have identical SSProve output/heap semantics.
\end{theorem}

The following theorem is the main program-logic bridge.

\begin{theorem}[Replace $q$ calls]
\label{thm:compile-replace}
Suppose two implementations \(P,P'\) of a selected operation satisfy a
one-call Pythagorean judgment with nonnegative tuple \(s\), preserve invariant
\(I\), and meet the compiler's typing, footprint, and memory-separation
conditions.  For every concrete SSProve program \(A\), replacing its first
\(q\) selected calls yields
\[
\begin{aligned}
 &\vDash\AEJ\{P_I\}\ \mathsf{Comp}_q(A;P,P)\\
 &\quad\approx_{\sqrt{q\norm{s}_1/2}}
   \mathsf{Comp}_q(A;P',P)\ \{o_L=o_R\},
\end{aligned}
\]
where \(P_I\) synchronizes inputs and heaps satisfying \(I\).
\end{theorem}
\begin{proof}[Proof sketch]
We first show that $\mathsf{Comp}_q(A;P,P)$ and $\mathsf{Comp}_q(A;P',P)$
satisfy the expected Pythagorean judgment (i.e. differ by the tuple $s^q$)
by induction on $q$ and applying our Pythagorean sequence rule.
The $q$-th call contributes \(s\) and preserves the invariant \(I\).

After \(q\) calls, the accumulated tuple has norm
\(q\norm{s}_1\); \textsc{Micciancio--Walter} is applied once, giving the
desired additive error bound.
\end{proof}

The theorem is generic in the program, packages, selected oracle,
invariant, and budget tuple.  The noise-flooding game supplies only the
one-call premise from \Cref{lem:one-call}; the program logic and compiler then
handle adversarial control flow.  Combined with
\Cref{thm:compile-same}, this is our quantitative counterpart of SSProve's
local relational reasoning for package indistinguishability.

\subsection{Why the extra judgment is necessary}

The additive-error layer remains essential for exact hops, up-to-bad
arguments, and the final triangle calculation.  Its ordinary sequence rule,
however, maps \(\delta_1,\delta_2\) to \(\delta_1+\delta_2\).  Applying Pinsker
inside each decrypt proof would therefore pay
\(q\sqrt{\epsilon_{\mathsf{nf}}/2}\).  The Pythagorean layer delays that
nonlinear conversion and obtains \(\sqrt{q\epsilon_{\mathsf{nf}}/2}\).  It is
not a new security definition; it is a verified relational abstraction that
preserves a parameter-critical cryptographic argument across adaptive,
stateful program composition.
